%%%%%%%%%%%%%%%%%%%%%%%%%%%%%%%%%%%%%%%%%
% University Assignment Title Page 
% LaTeX Template
% Version 1.0 (27/12/12)
%
% This template has been downloaded from:
% http://www.LaTeXTemplates.com
%
% Original author:
% WikiBooks (http://en.wikibooks.org/wiki/LaTeX/Title_Creation)
%
% License:
% CC BY-NC-SA 3.0 (http://creativecommons.org/licenses/by-nc-sa/3.0/)
% 
% Instructions for using this template:
% This title page is capable of being compiled as is. This is not useful for 
% including it in another document. To do this, you have two options: 
%
% 1) Copy/paste everything between \begin{document} and \end{document} 
% starting at \begin{titlepage} and paste this into another LaTeX file where you 
% want your title page.
% OR
% 2) Remove everything outside the \begin{titlepage} and \end{titlepage} and 
% move this file to the same directory as the LaTeX file you wish to add it to. 
% Then add \input{./title_page_1.tex} to your LaTeX file where you want your
% title page.
%
%%%%%%%%%%%%%%%%%%%%%%%%%%%%%%%%%%%%%%%%%
%\title{Title page with logo}
%----------------------------------------------------------------------------------------
%  PACKAGES AND OTHER DOCUMENT CONFIGURATIONS
%----------------------------------------------------------------------------------------
 
\documentclass[10pt]{article}
\usepackage[english]{babel}
\usepackage[utf8x]{inputenc}
\usepackage{amsmath}
\usepackage{graphicx}
\usepackage[colorinlistoftodos]{todonotes}
\usepackage{url}
\newcommand{\red}[1]{\textcolor{red}{#1} }
\usepackage{multirow}
\usepackage{booktabs} 
\usepackage{latexsym}
\usepackage{amssymb}
\usepackage{amsmath}
\usepackage{stmaryrd}
\usepackage{xspace}
\usepackage{url}
\usepackage{alltt}
\usepackage{multirow}
\usepackage{todo}

\newboolean{showcomments}
\setboolean{showcomments}{true} % toggle to show or hide comments
\ifthenelse{\boolean{showcomments}}
  {\newcommand{\nb}[2]{
    \fcolorbox{gray}{yellow}{\bfseries\sffamily\scriptsize#1}
    {$\blacktriangleright$#2$\blacktriangleleft$}
   }
   \newcommand{\version}{\emph{\scriptsize$-$working$-$}}
  }
  {\newcommand{\nb}[2]{}
   \newcommand{\version}{} 
  }
\newcommand\levi[1]{\nb{Levi}{\textcolor{blue}{#1}}}
\newcommand\vivek[1]{\nb{Vivek}{\textcolor{red}{#1}}}
\newcommand\ben[1]{\nb{Ben}{\textcolor{teal}{#1}}}


\begin{document} 

\begin{titlepage}

\newcommand{\HRule}{\rule{\linewidth}{0.5mm}} % Defines a new command for the horizontal lines, change thickness here

\center % Center everything on the page
 
%----------------------------------------------------------------------------------------
%  HEADING SECTIONS
%----------------------------------------------------------------------------------------

\textsc{\LARGE fortiss GmbH}\\[1.5cm] % Name of your university/college
\textsc{\Large Research Project}\\[0.5cm] % Major heading such as course name
% \textsc{\large Minor Heading}\\[0.5cm] % Minor heading such as course title

%----------------------------------------------------------------------------------------
%  TITLE SECTION
%----------------------------------------------------------------------------------------

\HRule \\[0.4cm]
{ \huge \bfseries A Case Study of Model Based Security Engineering for Industry 4.0}\\[0.4cm] % Title of your document
\HRule \\[1cm]
 
%----------------------------------------------------------------------------------------
%  AUTHOR SECTION
%----------------------------------------------------------------------------------------

\begin{minipage}{0.5\textwidth}
\begin{flushleft} \large
\emph{Authors:}\\
Vivek {Nigam} -- fortiss \qquad% Your name
\qquad Levi L\'ucio -- fortiss \qquad \qquad
\qquad Ben Schneider -- fortiss 
\end{flushleft}
\end{minipage}
~
\begin{minipage}{0.4\textwidth}
\end{minipage}\\[1cm]

% If you don't want a supervisor, uncomment the two lines below and remove the section above
%\Large \emph{Author:}\\
%John \textsc{Smith}\\[3cm] % Your name

%----------------------------------------------------------------------------------------
%  DATE SECTION
%----------------------------------------------------------------------------------------

{\large \today}\\[1.5cm] % Date, change the \today to a set date if you want to be precise

%----------------------------------------------------------------------------------------
%  LOGO SECTION
%----------------------------------------------------------------------------------------

\includegraphics[width = 0.3\textwidth]{figures/fortiss.png}\\[0.5cm] % Include a  
%----------------------------------------------------------------------------------------

\vfill % Fill the rest of the page with whitespace

\end{titlepage} 

\newcommand\autofocus{Auto{\sc focus}} 

\begin{abstract} 
Industry 4.0 has made huge progress by providing exciting new applications
facilitating, for example, factory management and configuration. However, less attention has been given on the security concerns introduced by the greater device connectivity. This is particularly bothersome with the use of technologies, such as, Internet of Things (IoT) which increases the attack surface of the system. These vulnerabilities have been witnessed by recent attacks involving IoT devices. The goal of this project is to study the use of Model Based Security Engineering techniques in building (more) secure Industry 4.0 systems. We will use the pick-and-place station already available at fortiss as main use-case modeling and analyzing its security concerns in an extension of \autofocus.
\end{abstract}

\section{Introduction} 
The new generation of manufacturing, the so-called Industry 4.0, is providing new and exciting features. Managers can easily configure an industry plant to the customer's needs, as well as help engineers to manage the factory by using, for example, plug-and-produce equipment. With the introduction of Internet of Thing (IoT) devices will provide even more features, such as accessing the factory remotely as well as allow the factory to use services on the cloud, such as IoT Watson. 

\textbf{While providing such features is desirable, such connectivity increases the attack surface of factories.} For example, if a compromised equipment is plugged in a factory, it \textbf{may gain access to secret resources or cause the factory to behave abnormally affecting the factory's productivity} or even cause some serious damage to the factory placing in danger the company's employees. In order to avoid such events, \textbf{it is important to place security counter measures}, such as the use of encryption keys, security protocols, and even physical barriers, \textbf{reducing thus vulnerabilities and security risk}. 

\begin{center}
  \emph{The main objective of this project is to apply Model-Based Security Engineering methodology for building secure Industry 4.0 applications.} 
\end{center}

\textbf{Model Based Security Engineering (MBSE)}~\cite{umlsec,secureuml} is a
methodology which brings together requirement and design engineers to carefully
analyze both safety and security concerns and their inter-dependencies. By
specifying requirements and verifying whether a design complies with given
requirements, \textbf{engineers can avoid attacks reducing security risks.} MBSE
can also help in determining \textbf{which counter-measures and where in a
system} should be placed in order to reduce security risks.
\begin{itemize}
  \item If a \textbf{security requirement specifies the integrity} of the data transmitted through a signal, then access control and signature mechanisms should be placed to enforce this requirement;

  \item In the same way, security requirements can specify whether some
  information flow from some (untrusted) signal should not be used by some
  critical device. This requirement shall trigger counter measures that avoid such information flows to occur;

  \item Models, such as Attack-Defense Trees, help engineers in
  \textbf{organizing threats, desired security properties, counter-measures and
  their (casual) relationships.} Moreover, by using annotations, such as the
  likelihood of threats and the damage caused by security property violation, it
  is possible to automate security risk analysis providing indications on where
  high security risks exist and where suitable counter-measures should be
  placed.
\end{itemize}
MBSE also guides the types of verification that can be deployed, for example, formal verification or security testing.

A substantial body of work on MBSE exists in the literature~\cite{advancesmds}.
UMLSec or SecureUML~\cite{umlsec,secureuml} are among the known approaches
to MBSE. Both use well-known UML diagrams such as \emph{Use Case diagrams},
\emph{Class diagrams}, \emph{Statecharts} or \emph{Sequence diagrams} to model
the business and the security aspects of an application. Security aspects can be embedded in the models either as orthogonal
languages that are composed with the business modeling languages, or by making
use of available facilities or annotation mechanisms already present in those
languages. Model transformations are then often used in MBSE to either produce
code or test- and verification-related artifacts directly from the models. This
allows elevating the level of abstraction at which the security aspects of an
application are modeled and verified.

\textbf{Automated model verification is one of the major advantages of MBSE.}
Most MBSE approaches use independent back-end verification technologies such as
model checkers or theorem provers. Such verification technologies are compatible 
with the abstraction levels at which security models are typically specified.

\section{Resources}

\subsection{Available Resources}

\subsection{Requested Resourced}

\paragraph{Available Resources and Requested Resources} 
This is an internal project in fortiss involving the Model-Based Software Engineering and the Industry 4.0 competence fields. 
\begin{itemize}
  \item \autofocus\ is a model-based framework developed in fortiss built on top
  of the FOCUS mathematical framework. It supports: requirements, including
  glossaries, data dictionaries and meta-data in the form of aspects; model elements for component design (code specifications, automata, etc); traceability of requirements and model elements; code generation; and integration with formal verification tools (Z3 and NuSMV). However, \textbf{\autofocus\  does not currently support Model-Based Security Engineering that allow the specification of security requirements, risk analysis models, and security model elements (crypto-languages, etc)};


  \item Pick-and-Place Component -- The Festo Didactic pick-and-place station is
  one of the factory elements maintained at fortiss. Its purpose is to place a
  cap on a cylinder. In particular, it contains all the necessary sensors to
  detect when the cylinder is at the correct place, and components to pick the
  cap from a station and place it on the cylinder. Moreover, the pick-and-place
  supports ``plug-and-produce'', allowing automatic connection and work with
  ``plug-and-produce'-enabled devices'. \textbf{However, no security
  counter measures are currently in place to defend against malicious agents.
  An example of such agents are malicious components that may gather information
  about the factory or even send incorrect events to disrupt the factory's production.}
\end{itemize}

\paragraph{Opportunity 1:} \textbf{The completion of this project will greatly
develop competence on MBSE for Industry 4.0}. On the one hand, \autofocus\
provides a platform for building prototypes supporting MBSE. On the other hand, the Pick-and-Place component provides an ideal example for applying MBSE;

\paragraph{Opportunity 2:} We expect to show the outcomes of this project at
fortiss' next Open House to be held after the 2018 summer break. \textbf{The
completion of this project will allow fortiss to advertise its competence on MBSE using the realistic (and eye catching) example of Industry 4.0.}

\paragraph{Opportunity 3:} We expect a close interaction between the Model-Based Software Engineering and Industry 4.0 competence fields. \textbf{This will lead to further collaborations between these (complementary) groups.}

\bigskip


The duration of this project is of 6 months. In order to carry out this project, we require the following resources:

\begin{itemize}
  \item 15\% Funding for Vivek Nigam;
  \item 15\% Funding for Levi L\'ucio;
  \item 15\% Funding for Ben Schneider;
  \item HiWi-1 of 20 hours per week;
  \item HiWi-2 of 20 hours per week.
\end{itemize}

\section{Tasks}

For this project we envision the following tasks:

\begin{itemize}
  \item Task 0: State of the art reviewing;
  \item Task 1: Vulnerability and Risk Analysis of the Pick-and-Place example;  
  \item Task 2: Security Requirements Engineering for the Pick-and-Place
  example;
  \item Task 3: Extension of \autofocus\ for Model Based Security Engineering
  support example;
  \item Task 4: Modeling of the Pick-and-Place example in \autofocus;
  \item Task 5: Investigation and Deployment of formal verification techniques for the Pick-and-Place Model;
  \item Task 6: Wrap-up with lessons learned.
  % \item Task 6: Translation of \autofocus\ design to 4Diac specification;
  % \item Task 7: Security Test-Case generation from the \autofocus\ model;
  % \item Task 8: Carrying out 
\end{itemize}

\paragraph{Task 0:} This project will start with an assessment of the current
state-of-the-art of MBSE. We will be particularly interested in identifying methodologies for writing security requirements, including identification of threats and counter-measures; and models for reasoning about security issues, such as Attack-Defense trees.

\paragraph{Task 1:} We will carry out a systematic vulnerability and risk
analysis of the Pick-and-Place component. This includes the identification of
desired security properties and the damage caused by their violation; entry
points and the likelihood that an attacker gains access to such entry points;
and counter measures to be used.

\paragraph{Task 2:} Based on the vulnerability analysis, we will develop a set
of security requirements. These requirements will be written in \autofocus\ and
will include enough meta-data to allow for the use of formal verification
techniques. More specifically we will develop new security aspects for
\autofocus, for example vulnerability, threats, assets and countermeasure
aspects.

\paragraph{Task 3:} Since \autofocus\ does not currently support MBSE, we will
extend it with features that will allow applying MBSE to the pick-and-place
component.
This includes adding support for security requirements; adding model elements
for risk analysis such as attack-defense trees; and extending \autofocus\
language with a crypto-language. These elements should be integrated with \autofocus'
requirements and aspects.

\paragraph{Task 4:} Given the identified security requirements and \autofocus\
extensions, we plan to model the Pick-and-Place example in \autofocus. This
includes adding traces from security requirements to model elements and
implementing identified counter measures.

\paragraph{Task 5:} We will also investigate the use of formal verification
techniques. This includes automating risk analysis and the use of
formal verification for identification of security flaws. For this we will use
verification tools that are already integrated with \autofocus, namely Z3 and
NuSMV. We shall extend existing \autofocus' model transformations into NuSMV
to include a crypto-language and security requirements expressed as aspects.

\paragraph{Task 6:} At the end of the project we will write a technical report
with the lessons learned. We also expect some publications from the work developed for
this project.


\subsection{Resource Allocation and Timetable}

We will divide the activities into two different teams. One team will be led by
Ben Schneider and will work on the tasks closer to the Industrial Automation
Group. The second team, led by Vivek Nigam, will work on the tasks
closer to the Model-Based Software Engineering group.

\begin{itemize}
  \item \textbf{Team 1:} Ben Schneider and one HiWi. This team will concentrate on the tasks more directly related to the Pick-and-Place, namely Tasks 1, 2 and 4;
  \item \textbf{Team 2:} Vivek Nigam, Levi Lúcio and one HiWi. This team will
  concentrate on the tasks related to the MBSE, namely, Task 3, 4, and 5.
\end{itemize}
Notice that tasks 0, 4 and 6 will be carried out by both teams.

Table~\ref{table:timetable} depicts the timetable of projects.  

\begin{table}[h]
\begin{center}
 \begin{tabular}{p{2cm}|p{1cm}|p{1cm}|p{1cm}|p{1cm}|p{1cm}|p{
1cm} }
\toprule
  \quad Month & \quad 1 &\quad 2 &\quad 3 &\quad 4 &\quad 5 &\quad
6\\
\midrule
  \quad Task 0 & \quad\checkmark &  & & & &\\
  \midrule
  \quad Task 1 & \quad\checkmark & \quad\checkmark & \quad\checkmark & & &\\
\midrule
  \quad Task 2 & &  & \quad\checkmark & \quad\checkmark & &\\
\midrule
  \quad Task 3 & \quad\checkmark &  \quad\checkmark & \quad\checkmark & \quad\checkmark & &\\
\midrule
  \quad Task 4 &  &  & &\quad\checkmark  & \quad\checkmark &\\
\midrule
  \quad Task 5 &  &  & & & \quad\checkmark& \quad\checkmark\\
\midrule
  \quad Task 6 &  &  & & & &\quad\checkmark\\
\bottomrule
 \end{tabular}
\end{center}
\caption{Timetable of the Project}
\label{table:timetable}
\end{table}


\section{Expected Results and Future Directions} 

We list some of the results expected from the completion of this project:

\begin{itemize}
  \item Security Requirements for the Pick-and-Place example;
  \item \autofocus\ Security Plug-In Development;
  \item \autofocus\ Model for the Pick-and-Place example;
  \item Identification of Security Risks and Possible Counter-Measures for the
  Pick-and-Place example;
  \item Closer interaction between Industry 4.0 and Model-Based Software Engineering groups;
  \item Competence development on Model Based Security Engineering for Industry 4.0;
  \item Presentation of our results at the fortiss Open-House 2018;
  \item Technical report and publications.
\end{itemize}


There are many possibilities that may emerge from this project. It should be
possible to apply the competence on MBSE and the machinery built in \autofocus
to other domains, such as Avionics, IoT applications, and other Cyber-Physical
systems. In the medium-term, the modeling of the Pick-and-Place will help us
understand the differences between \autofocus and 4diac, the framework
developed in the Industry 4.0 group. A possibility would then be to leverage the
modeling capabilities of \autofocus, such as connections with formal verification, and
the capabilities of 4diac, such as automatic code deployment.


\bibliographystyle{plain} 
\bibliography{bib}

\end{document}